\documentclass[11pt]{article}
\usepackage[margin=0.7in]{geometry}
\usepackage[T1]{fontenc}
\usepackage[utf8]{inputenc}
\usepackage{lmodern}
\usepackage{array}
\usepackage{booktabs}
\usepackage{enumitem}
\usepackage{hyperref}
\usepackage{longtable}
\usepackage{parskip}
\usepackage{ragged2e}
\usepackage{setspace}
\usepackage{tabularx}
\usepackage{titlesec}
\usepackage{xcolor}

\definecolor{PrimaryRed}{HTML}{8B1E1E}
\definecolor{SoftBlue}{HTML}{DDEEFF}
\definecolor{DarkGray}{HTML}{333333}

\hypersetup{colorlinks=true,linkcolor=black,urlcolor=black,citecolor=black}
\setlength{\parindent}{0pt}
\newcolumntype{Y}{>{\raggedright\arraybackslash}X}
\titleformat{\section}{\large\bfseries\color{PrimaryRed}}{}{0em}{}
\titleformat{\subsection}{\normalsize\bfseries}{\thesubsection}{0.5em}{}
\setstretch{1.12}

\begin{document}

\begin{titlepage}
\newgeometry{margin=0.85in}
\thispagestyle{empty}

\begin{center}
{\large University of Rochester\\}
{\normalsize Warner School of Education\\[0.35in]}

{\color{PrimaryRed}\rule{\textwidth}{1.4pt}}\\[0.22in]
{\Huge\bfseries Designing a Minecraft World\\[0.08in] That Carries Identity and\\[0.08in] Community Knowledge}\\[0.18in]
{\Large Class 7 Asynchronous Mini-Unit Draft}\\[0.22in]
{\color{PrimaryRed}\rule{0.78\textwidth}{0.7pt}}\\[0.32in]

\colorbox{SoftBlue}{%
\parbox{0.88\textwidth}{%
\centering\vspace{0.12in}
\textbf{Layered Genius: Designing for Joy and Identity}\\[0.05in]
Grade 4--5 Computer Science / Digital Literacy mini-unit centered on Minecraft Education, multimodal meaning-making, identity-safe design, and community knowledge
\vspace{0.12in}
}%
}
\end{center}

\vspace{0.38in}

\begin{center}
\renewcommand{\arraystretch}{1.35}
\begin{tabular}{>{\bfseries\RaggedLeft}p{1.65in} p{4.9in}}
Student & Piter Zacari Garcia Bautista \\
Course & EDU498: Literacy Learning as Social Practice \\
Class Session & Class 7: Towards a Pursuit of Identity \\
Assignment & Layered Genius: Designing for Joy and Identity \\
Mini-Unit Focus & Minecraft world-building as identity, community, and coding literacy \\
Student Profile & Ethan, Student 5, Computer Science \\
Submission Date & June 10, 2026 \\
\end{tabular}
\end{center}

\vspace{0.34in}

\begin{center}
\begin{minipage}{0.86\textwidth}
\itshape
This draft uses Piter's Puzzle Plan and EQUITAS lens to treat Minecraft Education as a literacy space where students can build, code, debug, explain, and represent meaning without being reduced to one narrow academic mode.

\vspace{0.22in}

\tiny\centering\textbf{Author's Context Note:} The draft is intentionally focused on the two pursuits named for this Class 7 activity, \textbf{Joy} and \textbf{Identity}. Skills, Intellect, and Criticality are preserved as brainstorm notes for the next revision so the lesson can grow across Muhammad's full five-pursuit framework without overloading the first submission.
\end{minipage}
\end{center}

\restoregeometry
\end{titlepage}
\clearpage
\onehalfspacing

\section*{Mini-Unit Snapshot}

\textbf{Class 7 task:} draft one lesson plan for a mini-unit that fully engages \textbf{Joy} and \textbf{Identity}, while leaving \textbf{Skills}, \textbf{Intellect}, and \textbf{Criticality} as brainstorm notes for later weeks.

\textbf{Grade / discipline:} Grade 4--5 Computer Science / Digital Literacy

\textbf{Student profile:} Ethan, a Seneca / Haudenosaunee and mixed-identity student who enjoys coding, game design, robotics, collaboration, outdoor learning, and community storytelling.

\textbf{Mini-unit focus:} Students use Minecraft Education to plan and begin building a small digital place that communicates identity, community knowledge, belonging, or future possibility.

\textbf{Essential question:} How can code, design, and digital place-making communicate who we are, what matters to us, and what our communities know?

\textbf{Three-lesson arc:}
\begin{itemize}[leftmargin=*]
    \item \textbf{Lesson 1: Plan a Meaningful Place.} Choose a place concept, sketch it, name meaningful features, and begin one prototype.
    \item \textbf{Lesson 2: Build and Debug.} Use Minecraft / MakeCode to build, test, revise, and add at least one coded or repeated design feature.
    \item \textbf{Lesson 3: Share and Reflect.} Present the world, explain design choices, and respond to peer feedback.
\end{itemize}

\section*{Source Trail and Rationale for the Pick}

This draft grows from the Class 7 prompt, the Class 7 tracker, the earlier base-selection memo, the Class 3 joy commitment, and the Grade 4 Minecraft placement evidence. The closest base was the existing Minecraft Education coding lesson family because it already included real evidence of debugging, peer teaching, block-counting, and multimodal digital literacy.

The other options were useful but less direct for this first draft. Squishy circuits would fit Sofia and embodied science literacy well, but the unit arc was less developed. AI/coding-as-literacy would be powerful for EQUITAS and criticality, but it was too complex for a first Joy + Identity lesson.

\section*{Student Profile and Identity Safety}

Ethan's profile matters because the lesson can recognize him as a coder, designer, storyteller, and community-minded problem solver. The lesson should not ask him to perform culture for the class. Instead, it offers identity-safe choice: students may build from lived experience, imagination, community, interests, family, land, language, hobbies, or future hopes.

Because Ethan's profile includes Seneca / Haudenosaunee heritage, culturally specific symbols or stories should be student-chosen and handled with care. Private disclosure is never required. Fictionalized, symbolic, or aspirational places are valid.

\section*{Pursuits Alignment}

\begin{tabularx}{\textwidth}{
>{\raggedright\arraybackslash}p{0.1\textwidth}
>{\raggedright\arraybackslash}p{0.9\textwidth}
}
\toprule
\textbf{Pursuit} & \textbf{How this draft addresses it} \\
\midrule
Joy & Students experience choice, visible progress, peer support, debugging as discovery, and the satisfaction of building something shareable. \\
Identity & Students design a meaningful place and explain how specific features communicate belonging, community, memory, or future possibility. \\
Skills & Brainstorm: sequencing, decomposition, loops, debugging, spatial reasoning, screenshot use, and oral/written explanation. \\
Intellect & Brainstorm: how code, design, space, and audience work together to create meaning. \\
Criticality & Brainstorm: whose stories appear in digital worlds, whose are missing, and how accessibility or respectful remixing shape design. \\
\bottomrule
\end{tabularx}

\section*{Multimodal and Multiliteracies Design}

Minecraft functions here as a text students author. They read digital space, sketch a plan, build visually, use procedural thinking, explain orally, label in writing, and revise through peer feedback.

\textbf{Modes included:} spatial design, visual design, coding/sequencing, oral explanation, written labels or reflection, screenshots, collaboration, and peer feedback.

\textbf{Access moves:} sketch-first or build-first choice; projected model; partner roles; sentence stems; optional oral explanation; screenshots as memory support; chunked directions; low-floor/high-ceiling challenge options.

\section*{Lesson 1 Plan}

\textbf{Lesson title:} Planning and Prototyping a Meaningful Minecraft Place

\textbf{Duration:} 60 minutes

\textbf{Student-friendly target:} I can plan and begin building a Minecraft place that shows something important about identity or community, and I can explain how my design choices carry meaning.

\textbf{Objectives:}
\begin{enumerate}[leftmargin=*]
    \item Students will choose a meaningful place connected to identity, community, belonging, or future possibility.
    \item Students will break the place into smaller build/code steps.
    \item Students will begin one prototype and test or revise at least one feature.
    \item Students will explain one design choice using meaning, audience, and purpose.
\end{enumerate}

\textbf{Materials:} Minecraft Education, teacher model world, projector, planning sheet, pencils or markers, optional MakeCode reference cards, and exit slip. If devices fail, students complete a paper-world storyboard and mark which steps would later be coded.

\textbf{Standards connection:} CSTA 1B-AP-10, 1B-AP-11, 1B-AP-15, and 1B-AP-17; ISTE Innovative Designer, Computational Thinker, Creative Communicator, and Global Collaborator.

\begin{tabularx}{\textwidth}{
>{\raggedright\arraybackslash}p{0.06\textwidth}
>{\raggedright\arraybackslash}p{0.64\textwidth}
>{\raggedright\arraybackslash}p{0.30\textwidth}
}
\toprule
\textbf{Mins} & \textbf{Activity} & \textbf{Purpose} \\
\midrule
0--7 & Show a generic Minecraft structure and a meaningful place. Ask: ``What makes a place meaningful, not just well built?'' & Frames design as meaning-making. \\
7--15 & Model a small place plan: entrance, sign, path, gathering spot, and one repeated pattern that could be built or coded. & Makes decomposition and design choices visible. \\
15--25 & Students choose a real, imagined, symbolic, or future place and sketch 3--5 important features. & Builds identity-safe agency. \\
25--38 & Students label 3--5 build/code steps and begin a prototype in Minecraft or on paper. & Moves from idea to manageable action. \\
38--50 & Students test one feature, debug one issue, or revise one design choice. & Treats mistakes as part of learning. \\
50--57 & Partner share: ``This place matters because...'' and ``One choice I made was...'' & Centers oral literacy and audience. \\
57--60 & Exit slip: ``What does your place show? What is your next build/code step?'' & Captures evidence and next steps. \\
\bottomrule
\end{tabularx}

\textbf{Differentiation:} {Students may work alone or with roles such as builder, navigator, debugger, or explainer. Beginners may build one strong scene; advanced students may add loops, agent movement, signage, or an audience-centered accessibility feature.}

\section*{Assessment}

\textbf{Evidence collected:} planning sheet, teacher conference notes, prototype or screenshot, partner explanation, and exit slip.

\textbf{Success criteria:}
\begin{enumerate}[leftmargin=*]
    \item The place has a clear connection to identity, community, belonging, or future possibility.
    \item The student breaks the idea into manageable build/code steps.
    \item The student begins, tests, or revises one prototype feature.
    \item The student explains how at least one design choice carries meaning.
\end{enumerate}

\section*{Reflection and Next Steps}

This draft keeps the assignment focused. Joy is built into choice, creation, debugging, peer support, and visible success. Identity is built into meaningful place-making and student control over what is shared. The lesson also protects privacy, which matters because identity-centered teaching should not require students to disclose trauma, culture, family history, or community knowledge on demand.

Muhammad's identity pursuit helps frame the work as self-definition, not decoration. Moje, Luke, Davies, and Street help ask what identity the task makes possible. Skerrett supports the idea that students can be repositioned as capable contributors. Moje's disciplinary literacy frame helps name coding, debugging, design, and explanation as literacy practices.

% The next revision should expand Lesson 2 and Lesson 3 so Skills, Intellect, and Criticality become more explicit without crowding the first lesson.


\newpage
\section*{References and Working Sources}

\begin{itemize}[leftmargin=*]
    \item EDU498 Class 7 prompt, \emph{Towards a Pursuit of Identity}: ``Layered Genius: Designing for Joy and Identity.''
    \item Muhammad, G. (2020). \emph{Cultivating Genius}, Chapter 3: Identity.
    \item Moje, E. B., Luke, A., Davies, B., \& Street, B. (2009). ``Literacy and identity.''
    \item Skerrett, A. (2012). ``We hatched in this class.''
    \item Moje, E. B. (2015). ``Doing and teaching disciplinary literacy with adolescent learners.''
    \item Class 3 joy commitment: ``belonging before building.''
    \item EDU498 Grade 4 Minecraft placement evidence: debugging, collaborative counting, peer teaching, and student-led problem framing.
    \item CSTA K--12 Computer Science Standards; ISTE Standards for Students; Minecraft Education / MakeCode support materials.
\end{itemize}

\end{document}